\documentclass[main.tex]{subfiles}
\begin{document}

In the U.S., Visa and Mastercard (MC) process 80\% of card transactions
% CONSTANT: from Visa 2020 annual report, external data
and earn profit margins above 60\% \parencite{Visa2020}.
% CONSTANT: from Visa 2020 annual report, external data
At the same time, merchants pay around \$120 billion in fees every year to accept cards \parencite{Nilson2020a}.
% CONSTANT: from Nilson (2020a), aggregate merchant fee bill
These facts have motivated two decades of litigation and legislation built on a view that high merchant fees reflect weak competition.
Key policy interventions include the Durbin Amendment's cap on debit interchange fees, the proposed Credit Card Competition Act \parencite{Andriotis2022}, and the Department of Justice (DOJ) monopolization lawsuit against Visa \parencite{Au-Yeung2024}.

I find that merchant fees are indeed too high, but weak competition is the wrong diagnosis.
Networks set rewards for consumers and fees for merchants. 
The division of costs matters because merchants typically charge the same price for all payment methods, a phenomenon known as ``price coherence'' \parencite{Frankel1998, Stavins2018}.\footnote{
    This occurs even though cash discounts and card surcharges are largely legal.
    I explore surcharging both theoretically and empirically in Online Appendix \ref{sec:oa-price-coherence}.
}
Under price coherence, card users capture rewards while the cost of higher merchant fees is spread across all consumers through higher retail prices.
I estimate that consumers are ten times more sensitive to rewards than merchants are to fees.
% HARDCODED[derived: diversion_baseline.tex "Visa Credit"/"Visa Credit" ÷ scalar divert_inst_merchant_price_visa_baseline] current=3.7/0.4=9.25
This incentivizes networks to charge high merchant fees to fund generous consumer rewards.
High-merchant-fee credit cards thus proliferate because of, not in spite of, intense network competition.
% Network driving adoption of high-fee payment methods and inflating merchants' costs.

Three forces explain why networks tax merchants to subsidize consumers.
First, consumer payment choices respond strongly to subsidies.
Networks charge interchange fees to merchants and channel the fees to card issuers, which incentivizes them to promote the networks' cards.
A regulatory reduction in debit interchange fees --- the Durbin Amendment --- reduced debit card spending at regulated issuers by around 25\% relative to exempt issuers.
% HARDCODED[table: durbin_table.tex, row: signature debit, exp(-0.287)-1, rounded] current=25
Second, card acceptance raises merchant sales.
Event-study evidence shows that merchants on the margin gain \scalarpctinput{kilts_model_sales_event}\% in sales when accepting credit.
% HARDCODED[scalar: kilts_model_sales_event (pct)] current=12
Third, while almost all merchants accept every network's cards, around \scalarinput{kilts_model_singlehome_cc}\% of consumers carry credit cards from only one network (i.e., they ``single-home'').
Together, these facts produce a partial ``competitive bottleneck'' \parencite{Armstrong2006}. 
Merchants risk losing substantial sales from single-homing consumers by dropping a network, so networks compete through rewards rather than fees.

Existing theory cannot resolve whether network competition raises or lowers fees and welfare.
The U.S.\ market sits between two theoretical extremes.
When all consumers single-home, each network is a monopoly gateway to its cardholders, so competition raises per-transaction merchant fees \parencite{Armstrong2006, Edelman.Wright2015}. 
When consumers multi-home (i.e. carry cards from multiple networks), merchants can drop expensive networks without losing sales, so competition lowers fees \parencite{Anderson.etal2018, Bakos.Halaburda2020, Teh.etal2022}.

I thus build a structural model in which payment networks compete as two-sided platforms.
A mix of single- and multi-homing consumers choose payment methods based on rewards, merchant acceptance, and exogenous non-price characteristics.
Merchants choose which cards to accept and set retail prices to maximize profits net of merchant fees.
Multiproduct networks maximize profits by setting fees and rewards, balancing merchant acceptance against consumer adoption.

I estimate consumer and merchant preferences by matching the reduced-form facts, moments from payment surveys, and aggregate market-level data.
The estimated parameters match untargeted moments, including the acceptance response to American Express (AmEx) fee cuts, merchant profit margins, and network costs from accounting data.
A 1-bp increase in Visa credit rewards raises Visa's consumer market share by $3.7\%$, while a 1-bp increase in Visa credit merchant fees reduces merchant acceptance by only $0.4\%$.
% HARDCODED[scalar: divert_inst_merchant_price_visa] current=-0.4
% HARDCODED[table: diversion_baseline.tex, row: "Visa Credit", col: "Visa Credit"] current=3.7
Networks thus face strong incentives to raise merchant fees to fund rewards.

Credit card adoption is socially excessive because of price coherence.
Because merchants do not surcharge, each consumer who adopts a credit card raises retail prices for all other consumers.
Consumers individually face incentives to distort their payment choices to capture cross-subsidies, but collectively prefer lower credit card use.
A revealed preference argument illustrates the inefficiency.
Since credit cards pay rewards whereas cash and debit do not, the marginal user of credit cards must prefer the non-price characteristics of cash and debit.
High rewards reduce consumers' aggregate non-price utility from payment choices, but the rewards themselves are merely transfers funded by higher retail prices.
% Credit fee caps correct this prisoner's dilemma by reducing the rewards that drive excessive adoption.

The structural model measures the welfare loss from this coordination failure.
I simulate a counterfactual in which credit card merchant fees are capped at 120 basis points (bp), roughly half their current level.
% CONSTANT: counterfactual design choice, approximately half current credit card fees
% \footnote{Online Appendix~\ref{subsec:oa-cost-of-cash-cap} evaluates alternative fee caps.
% The 120 bps cap achieves roughly 93\% of the welfare gains from a 30 bps cap modeled after the tourist test of \textcite{Rochet.Tirole2011} and roughly 77\% of a Ramsey planner's gains.
% % HARDCODED[table: appendix_cap_welfare_table_baseline.tex, row: "Total", cols: Cap Credit=27, Tourist Test=29, Ramsey Planner=35] current=93% (27/29), 77% (27/35)
% }
Capping merchant fees cuts per-transaction revenue, so networks compete less aggressively for consumers through rewards.
Networks cut rewards, and some consumers switch away from credit cards.
By revealed preference, these marginal switchers prefer cash or debit absent rewards, so the total welfare gain is the eliminated cross-subsidy.
This reduction in credit card use ultimately raises total welfare by \$27 billion annually.
The gains from capping merchant fees exceed the \$12 billion consumer welfare gain from the CARD Act \parencite{Agarwal.etal2015},
% CONSTANT: from Agarwal et al. (2015), external paper result
suggesting that regulating networks is at least as important as regulating issuers.

The gains from credit fee caps are also progressive. 
Reduced merchant fees pass through to lower retail prices for all consumers, whereas the reduction in rewards falls mostly on high-income consumers who are more likely to use credit cards.
Consumer surplus rises by 48 bp of consumption for low-income consumers, compared to only 15 bp for high-income consumers.
% HARDCODED[table: welfare_table_baseline.tex, row: "Low Income", col: 1] current=48
% HARDCODED[table: welfare_table_baseline.tex, row: "High Income", col: 1] current=15
There is no trade-off between equity and efficiency in regulating payment fees.

Two additional counterfactuals illustrate that the largest welfare losses in payment markets stem from the overuse of credit cards, not market power.
First, repealing the Durbin Amendment's debit fee cap would raise total welfare by \$7 billion per year.
% HARDCODED[table: welfare_table_baseline.tex, row: "Total", col: 3] current=7
Under the usual market power story \parencite{Cuesta.Sepulveda2021}, removing the debit fee cap should reduce welfare.
But higher debit merchant fees fund higher debit rewards, drawing consumers away from high-reward credit cards onto a lower-fee network.
This shrinks the cross-subsidy and increases total welfare.
This substitution effect matches prior work on the Durbin Amendment \parencite{Mukharlyamov.Sarin2025} and illustrates that the current U.S. regime of capping debit but not credit fees is worse than laissez-faire.

Second, a merger to monopoly raises total welfare by \$15 billion per year.
% HARDCODED[table: welfare_table_baseline.tex, row: "Total", col: 6] current=15
Without competitive pressure to attract consumers, a monopoly network cuts rewards.
Because consumers are highly sensitive to rewards, credit card use collapses.
Per-transaction merchant fees rise, but credit card volume collapses, so the total merchant fee burden falls.
Consumers lose \$11 billion as network markups increase sharply, but the reduction in credit card use raises total welfare.
% HARDCODED[table: welfare_table_baseline.tex, row: "Consumers", col: 6] current=-11
While a merger is not a realistic policy proposal, the counterfactual shows how network competition inflates the total merchant fee burden by funding rewards.
% HARDCODED[table: welfare_table_baseline.tex, row: "Total", col: 6] current=15
These predictions align with rising U.S.\ interchange fees despite increasing network competition \parencite{GAO2009} and the global expansion of high-fee Buy Now, Pay Later products \parencite{Berg.etal2024}.
The problem is not market power stifling output but rewards competition generating excessive credit card adoption.

In the absence of credit fee caps, policy can also redirect competition by increasing consumer multi-homing.
Dual-routing requirements, such as those proposed in the Credit Card Competition Act, mandate that issuers put multiple networks on each card.
Cardholders covered by such mandates multi-home by technological fiat, without having to sign up for a second card.
I model these requirements as a change in preferences that generates a 20 percentage point increase in the share of multi-homing consumers.
When consumers multi-home, merchants can decline expensive networks without losing sales, shifting competition toward merchant fees rather than consumer rewards.
Networks respond by cutting credit card fees by 13 bp and rewards by 15 bp.
% HARDCODED[table: cf_table_baseline.tex, row: "Credit Fees", col: 8] current=-13
% HARDCODED[table: cf_table_baseline.tex, row: "Credit Rewards", col: 8] current=-15
Total welfare rises.
These routing requirements are effective even when the secondary network is a large player.
Increasing multi-homing is also much more effective than introducing a public option, such as a Central Bank Digital Currency (CBDC), that competes as a debit card \parencite{Berg.etal2024a}.
High rates of consumer multi-homing, not low market shares, force networks to compete for merchants rather than consumers.

\paragraph{Related Literature}

The closest empirical work is \textcite*{Huynh.etal2022}, which estimates a structural two-sided model of consumer and merchant card adoption.
I build on their work by modeling merchant and network competition.
Merchant competition explains how credit card rewards inflate retail prices and redistribute consumption.
Network competition endogenizes merchant fees and consumer rewards, enabling me to assess how price controls and competition affect total welfare.

The effects of network competition on fees and welfare are quantitative questions, but common modeling assumptions predetermine the answer.
Homogeneous merchants guarantee that competition reduces fees \parencite{Guthrie.Wright2007}.
Omitting merchant competition understates merchants' incentives to accept cards \parencite{Rochet.Tirole2002, Wright2012, Li.etal2020} and misses how fees pass through to retail prices.
My model relaxes both assumptions and lets the data determine whether competition raises or lowers fees and welfare.

My paper contributes to the empirical IO literature on two-sided markets \parencite{Rysman2004, Lee2013, Song2021}.
My results on dual routing support the finding in \textcite{Rosaia2020} that consumer multi-homing intensifies platform competition.
My paper also contributes to the literature on how platforms shape off-platform outcomes \parencite{Bergemann.etal2024, Chen2024, Bursztyn.etal2025}.
Under price coherence, merchant fees strongly influence retail prices, explaining the welfare gains from credit fee caps.
\textcite{Sullivan2025} finds that commission caps reduce welfare in food delivery, where price coherence is absent.

My paper also relates to a literature on price regulation in credit card markets.
Whereas \textcite{Nelson2020,Cuesta.Sepulveda2021} show how interest rate caps increase welfare by expanding quantities, my paper shows that credit card merchant fee caps can increase welfare by reducing the use of credit cards.
\textcite{Knittel.Stango2003} document that non-binding state interest rate caps served as focal points for tacit coordination in credit card markets during the 1980s.
The credit fee caps I study are binding constraints that push fees below their competitive levels, so they function as direct price regulation rather than as coordination devices.

\end{document}